\section{模型假设}

\begin{enumerate}
  \item \textbf{功率分段恒定。}每个 10 分钟区间内电价、负载与光伏功率视为常数，区间电量等于功率乘以 $1/6\,\mathrm{h}$。
  \item \textbf{充放电效率全记在充电侧。}按附录 1 的“充放电效率 $90\%$”，本文取充入 $1\,\mathrm{kWh}$ 有 $0.9\,\mathrm{kWh}$ 进入电池、放出 $1\,\mathrm{kWh}$ 电池减少 $1\,\mathrm{kWh}$，完整充放一次为 $0.9$。
  \item \textbf{不允许向外网售电。}光伏富余时只能给储能充电或弃光，因此计划中显式设置弃光变量 $w_t\ge0$，且购电量 $x_t\ge0$。
  \item \textbf{储能在日前计划内严格按曲线充放电。}执行阶段实际负载与光伏一旦偏离预测，差额全部由紧急购电（不足）或弃光（有余）吸收，储能不临时改变充放电量。这既是题面“制定计划购电策略”的自然含义，也是式~\eqref{eq:identity} 成立的前提。
  \item \textbf{日前决策严格因果。}第 $d$ 天 0:00 的计划只允许使用第 $d-1$ 天及更早的历史数据；问题三中 $6{:}00$、$12{:}00$、$18{:}00$ 的调整只允许使用该时刻已发布的预报。任何时刻都不得读取当天未来的实际值。
  \item \textbf{储能日循环运行。}问题一中储能 0:00 与 24:00 的储电量相等且均为 $6000\,\mathrm{kWh}$；问题二至四把该口径推广为每天计划结束时回到当天起始电量。由于电价日型固定，跨日套利的空间很小，该假设的代价在灵敏度分析中量化。
  \item \textbf{计划内弃光等价于预留裕度。}计划中出现的弃光量 $w_t>0$ 表示“按预测多买、多买的电在预测兑现时弃掉”，它是应对预测误差的保险，而非设备故障。
\end{enumerate}
